\chapter{Conclusion}\label{chap:conclusion}

This thesis set out to understand and model how news propagates into the liquidity of WTI crude oil futures at hourly resolution, along two axes: the temporal lag of the response (RQ1) and its directional asymmetry (RQ2). It went at the questions in two phases, a regression-based Phase~1 on FinBERT sentiment and a deep-learning Phase~2 on LLM-extracted, channel-decomposed features.

This chapter states the contributions, summarizes the findings, and closes.

\section{Summary of contributions}\label{sec:contributions-summary}

The thesis makes three empirical and five methodological contributions.

Empirically:

\begin{itemize}
  \item \textbf{An hourly lag structure for the news-liquidity response.} The response peaks at $+6$ hours, having built over the preceding hours and decaying to insignificance by $+8$. The Phase~1 lag OLS establishes it, and the Phase~2 TFT corroborates it independently through its direction-conditioned attention, which places the bearish peak at $-6$ hours under a method sharing no assumptions with the regression. Daily-frequency work cannot see this structure at all.
  \item \textbf{A precise account of the directional asymmetry.} Bearish sentiment carries a larger marginal coefficient than bullish at every dominant lag, by about 18\% at the peak (Phase~1). The ordering is consistent across lags, though the difference between the two coefficients is not itself submitted to an equality test (Section~\ref{sec:future-work}). It does not appear at all as a difference in the level of predicted volume (Phase~2), because the model organizes its predictions around risk and salience rather than bullish or bearish price direction. The precision here is as much the contribution as the result: stating the asymmetry correctly means naming the sentiment construct and the estimand it refers to.
  \item \textbf{A news-representation finding.} Title-only and title-plus-body FinBERT sentiment disagree on 41.6\% of articles, and titles lean more bullish than the bodies they head. Which news representation you pick is itself a modeling decision.
\end{itemize}

Methodologically (developed in Section~\ref{sec:methodological-contributions}):

\begin{itemize}
  \item \textbf{Temporal alignment of continuous news onto a discontinuous trading grid}, the prerequisite for any intraday design: a causally conservative ceiling rule plus forward-assignment of the 22.9\% of articles that publish while the market is shut.
  \item \textbf{Channel decomposition} as a higher-resolution feature set and, through cross-model agreement, a moderately defensible validation substrate where no human ground truth exists.
  \item \textbf{The tone-versus-price distinction}, showing that a sentiment measure is defined by its construct, and that a tone classifier and a price-reasoning LLM part ways on exactly the geopolitical news that drives volume.
  \item \textbf{LLM-as-filter} as a semantically grounded alternative to regex heuristics for noisy web-scraped news corpora.
  \item \textbf{A two-phase template} pairing an interpretable statistical baseline with an expressive deep-learning model that corroborates it rather than replacing it.
\end{itemize}

\section{Summary of findings}\label{sec:findings-summary}

On RQ1, news hits WTI liquidity hardest not in the same hour but roughly six hours later. Two independent methods agree on this: the Phase~1 lag OLS peaks at $+6$ hours, and the Phase~2 TFT, trained and evaluated separately, attends to $-6$ hours for bearish news. Its persistence-relative gain rises across the whole 1-to-12 hour window, confirming that the Phase~2 features carry signal throughout it, but that curve is driven by the decay of the baseline and does not itself single out a lag. The liquidity response to news is gradual rather than instantaneous.

On RQ2, the answer depends on how the question is asked. Bearish sentiment carries a consistently larger marginal coefficient than bullish across the dominant lags, an asymmetry in the sensitivity of volume to negative-tone news that is the thesis' primary answer, and one stated as a consistent ordering rather than a tested contrast. It does not survive as a difference in the average level of predicted volume, for two reasons: the model organizes its forecasts around the risk and salience a news item carries rather than its price direction, and a tone-based and a price-based sentiment measure disagree on exactly the geopolitical news that moves the market. WTI liquidity responds to how much a piece of news raises supply uncertainty more than to whether it reads as nominally good or bad.

As a forecasting exercise, TFT v2 cut \texttt{log\_volume} error over a persistence baseline by 46 to 71 percent across horizons and \texttt{amihud} by 43 to 45 percent, with VIX, the supply and demand channels, and the entity flags as a block, which carry 52 percent of total feature importance between them, as what it leans on. It failed only on \texttt{price\_range} in the unseen war regime, which is a clean case of regime extrapolation. Alongside the two headline answers, three supporting findings shaped how the research developed: a 41.6 percent title-versus-body sentiment disagreement, a composite-sentiment calibration failure between model families that motivated the channel decomposition, and a semantic-versus-lexical filter comparison that motivated the \texttt{usable\_strict} rule.

\section{Closing remarks}\label{sec:closing}

This thesis started from a simple observation, that news moves markets, and asked something more specific of it: at what lag, and with what asymmetry, does news move the liquidity of WTI crude oil at hourly resolution. The answer has two parts. Empirically, the response is gradual and peaks about six hours after publication, and its directional asymmetry is consistent but has to be stated carefully, because it lives in the sensitivity of volume to news rather than in the average volume of nominally good or bad hours. Methodologically, getting to that precision meant treating sentiment not as a single quantity handed to us but as a construct to define, decompose, and validate. That is what the channel decomposition, the LLM extraction, and the two-phase design are for.

If one idea carries beyond this study, it is that for a commodity like crude oil the liquidity response to news is organized around risk rather than direction. The market does not trade more because the news is bad. It trades more because the news raises uncertainty about supply, and in the oil market that uncertainty is often, against intuition, bullish for price. Once that is clear, the apparent tension between a tone-based and a price-based reading of the same events goes away, and it changes what a news-driven liquidity model should be built to detect.

The work is bounded: one commodity, two years, one structural break, and a validation resting on model agreement rather than human judgment. Within those bounds it offers a concrete empirical result and a reusable way of getting there, and the extensions in Section~\ref{sec:future-work} are mostly small steps rather than new programs. How news propagates into market liquidity is a question of resolution as much as of theory, and at hourly resolution it shows more structure, and more nuance, than a daily view can.
